\section{The uniform law and the $2\sin^2$ lemma}\label{cg:sec:uniform}

The shifted family does not lie in $\CGLP$ (Section~\ref{cg:sec:cmfamily}). The $2\sin^2$ lemma asks a different question: given a L\'evy measure on the primal side, what is the Thorin measure on the dual side?

\subsection{The lemma and the Thorin density of a signed L\'evy measure}

The $2\sin^2$ lemma is \eqref{eq:sin2} (Lemma~\ref{rz:lem:P1}, Theorem~\ref{gl:thm:sinsq}), proved there. In the form used here, for $x>0$ and $w\ge0$,
\begin{equation}\label{cg:eq:twosiny}
\frac{e^{-wx}+wx-1}{x}=\frac1\pi\int_0^\infty\log\Bigl(1+\frac{w^2}{y^2}\Bigr)\bigl(1-\cos xy\bigr)\,dy .
\end{equation}
In the variables $u=w^2/2$, $z=y^2/2$ this is $e^{-wx}+wx-1=\int_0^\infty\log(1+u/z)\,\tau_x(z)\,dz$ with $\tau_x(z)=\frac{x}{\pi\sqrt{2z}}\,2\sin^2\bigl(x\sqrt{z/2}\bigr)\ge0$.

The next statement differs from \eqref{gl:eq:sinsq-thorin} (Chapter~\ref{chap:gl}) only in its hypothesis, an abstract signed measure $M$ in place of the L\'evy measure of $\xi$.

\begin{theorem}[Thorin density from a signed L\'evy measure]\label{cg:thm:thorinfromlevy}
Let $f$ be analytic and zero free on $\operatorname{Re}s>\alpha_0$ and positive on $(\alpha_0,\infty)$. Let $M$ be a signed measure on $(0,\infty)$ with $\int(1\wedge x)e^{-\alpha x}|M|(dx)<\infty$ for every $\alpha>\alpha_0$, such that
\begin{equation}\label{cg:eq:levyk}
\frac{f'}{f}(s)-\frac{f'}{f}(\alpha)=\int_0^\infty\bigl(e^{-\alpha x}-e^{-sx}\bigr)M(dx),\qquad \operatorname{Re}s>\alpha_0,\ \alpha>\alpha_0 .
\end{equation}
Then for every $\alpha>\alpha_0$ and $w\ge0$,
\begin{equation}\label{cg:eq:thorinrep}
\log\frac{f(\alpha+w)}{f(\alpha)}-w\,\frac{f'}{f}(\alpha)=\int_0^\infty\log\Bigl(1+\frac{w^2}{y^2}\Bigr)\kappa_\alpha(y)\,dy,
\end{equation}
\begin{equation}\label{cg:eq:kappa}
\kappa_\alpha(y)=\frac1\pi\int_0^\infty\bigl(1-\cos xy\bigr)e^{-\alpha x}M(dx)=\frac1\pi\Bigl[\operatorname{Re}\frac{f'}{f}(\alpha+iy)-\frac{f'}{f}(\alpha)\Bigr].
\end{equation}
If $\kappa_\alpha\ge0$, then $u\mapsto F_\alpha(u)=\dfrac{f(\alpha)}{f(\alpha+\sqrt{2u})}\exp\bigl\{\sqrt{2u}\,\frac{f'}{f}(\alpha)\bigr\}$ is the Laplace transform of a GGC with Thorin density $\tau_\alpha(z)=\kappa_\alpha(\sqrt{2z})/\sqrt{2z}$.
\end{theorem}

\begin{proof}
Differentiating \eqref{cg:eq:levyk} in $s$ gives $(f'/f)'(s)=\int xe^{-sx}M(dx)$. The function $\int_0^\infty e^{-\alpha x}(e^{-wx}+wx-1)\,M(dx)/x$ has the same second derivative in $w$ as the left side of \eqref{cg:eq:thorinrep}, and both vanish together with their first derivatives at $w=0$, so they are equal. Now insert \eqref{cg:eq:twosiny}. Fubini is justified by Tonelli applied to $|M|$: by \eqref{cg:eq:twosiny} the double integral of the absolute values is $\int e^{-\alpha x}(e^{-wx}+wx-1)|M|(dx)/x\le\int e^{-\alpha x}\min(w^2x/2,w)\,|M|(dx)<\infty$. The second form of $\kappa_\alpha$ is \eqref{cg:eq:levyk} at $s=\alpha+iy$, since $\operatorname{Re}(e^{-\alpha x}-e^{-(\alpha+iy)x})=e^{-\alpha x}(1-\cos xy)$. If $\kappa_\alpha\ge0$ then $\log F_\alpha(u)=-\int\log(1+u/z)\tau_\alpha(z)\,dz$ with a nonnegative density and a finite integral, which is the definition of a GGC transform \cite{Bondesson}.
\end{proof}

When $M\ge0$, $\kappa_\alpha\ge0$ automatically. Formula \eqref{cg:eq:kappa} gives a test that needs no L\'evy measure: $\kappa_\alpha\ge0$ if and only if the real part of the logarithmic derivative on the vertical line through $\alpha$ is never below its value on the real axis.

\begin{corollary}[Zeros on a line give Lorentzians]\label{cg:cor:lorentz}
Let $a$ be real and let the hypotheses of Theorem~\ref{cg:thm:thorinfromlevy} hold. Suppose $f'/f(s)=a+\sum_k(s-\rho_k)^{-1}$ (symmetric summation) with all $\rho_k=\beta_0+it_k$ on one vertical line, the set $\{t_k\}$ symmetric and $0\notin\{t_k\}$. Then for $\alpha=\beta_0+\varepsilon$, $\varepsilon>0$,
\begin{equation}\label{cg:eq:lorentz}
\kappa_\alpha(y)=\frac1\pi\sum_k\bigl[\ell_\varepsilon(y-t_k)-\ell_\varepsilon(t_k)\bigr],
\end{equation}
and if $\sum_kt_k^{-2}<\infty$ then $\kappa_\alpha(y)\,dy\to\sum_{t_k>0}\delta_{t_k}(dy)$ vaguely on $(0,\infty)$ as $\varepsilon\downarrow0$; the right side of \eqref{cg:eq:thorinrep} tends to $\sum_{t_k>0}\log(1+w^2/t_k^2)$.
\end{corollary}

\begin{proof}
$\operatorname{Re}(\alpha+iy-\beta_0-it)^{-1}=\ell_\varepsilon(y-t)$. As $\varepsilon\downarrow0$, $\pi^{-1}\ell_\varepsilon(\cdot-t)\to\delta_t$; the constant $\pi^{-1}\sum_k\ell_\varepsilon(t_k)\le\pi^{-1}\varepsilon\sum_kt_k^{-2}\to0$; and peaks at $t_k<0$ do not reach $y>0$. For the last limit the integral against $\log(1+w^2/y^2)$ needs the tail bound $\sum_{k}O(w^2/t_k^2)$ for the dominated convergence, which the hypothesis $\sum t_k^{-2}<\infty$ supplies.
\end{proof}

\subsection{The tilted uniform law}

Take $f=\varphi_U(s)=(1-e^{-s})/s$. By Proposition~\ref{cg:prop:signed} (after the tilt, where the integrals converge absolutely), \eqref{cg:eq:levyk} holds with $M_U(dx)=dx-d\lfloor x\rfloor$, $f'/f=g$ and $\alpha_0=0$, because $g(s)=\sum_{n\ge1}e^{-ns}-1/s=\int_0^\infty e^{-sx}(d\lfloor x\rfloor-dx)$ and the zeros of $\varphi_U$ are $2\pi ik$, $k\ne0$, on the line $\operatorname{Re}s=0$. The tilt is $\alpha=\delta$: it replaces $g(x)$ by $g(x+\delta)$ in the Mellin integral, which is the shifted family of Section~\ref{cg:sec:cmfamily}.

\begin{theorem}[Thorin density of the tilted uniform law]\label{cg:thm:uniform}
For every $\delta>0$,
\begin{equation}\label{cg:eq:kappaU}
\begin{split}\pi\kappa_\delta(y)&=\frac{y^2}{\delta(\delta^2+y^2)}-\frac12\coth\frac\delta2\;\frac{1-\cos y}{\cosh\delta-\cos y}\\
&=\sum_{m\ge1}\bigl[\ell_\delta(y-2\pi m)+\ell_\delta(y+2\pi m)\\
&\qquad-2\ell_\delta(2\pi m)\bigr],
\end{split}\end{equation}
and $\kappa_\delta(y)>0$ for all $y>0$. Hence $u\mapsto\dfrac{\varphi_U(\delta)}{\varphi_U(\delta+\sqrt{2u})}\,e^{\sqrt{2u}\,g(\delta)}$ is the Laplace transform of a GGC with Thorin density $\tau_\delta(z)=\kappa_\delta(\sqrt{2z})/\sqrt{2z}$. As $\delta\downarrow0$, $\kappa_\delta(y)\,dy\to\sum_{m\ge1}\delta_{2\pi m}$ vaguely, and the limit of \eqref{cg:eq:thorinrep} is
\[
\log\varphi_U(w)+\frac w2=\log\frac{\sinh(w/2)}{w/2}=\sum_{m\ge1}\log\Bigl(1+\frac{w^2}{4\pi^2m^2}\Bigr),
\]
so $(w/2)/\sinh(w/2)$ is the Laplace transform in $u=w^2/2$ of $\sum_me_m/(2\pi^2m^2)$, a GGC with Thorin measure $\sum_m\delta_{2\pi^2m^2}$ (a law of Biane, Pitman and Yor type \cite{BPY01}; it is $S_1/4$ in their notation).
\end{theorem}

\begin{proof}
The closed form follows from \eqref{cg:eq:kappa} with $g(s)=\sum_{n\ge1}e^{-ns}-1/s$: one has $1/\delta-\operatorname{Re}(\delta+iy)^{-1}=y^2/\{\delta(\delta^2+y^2)\}$ and $\sum_ne^{-n\delta}(1-\cos ny)=\frac12\coth(\delta/2)(1-\cos y)/(\cosh\delta-\cos y)$. The Lorentzian form is Corollary~\ref{cg:cor:lorentz}, since $g(s)=-\frac12+\sum_{k\ne0}(s-2\pi ik)^{-1}$. For positivity let $P(y)=\sum_{m\in\mathbb Z}\ell_\delta(y-2\pi m)=\frac12\sinh\delta/(\cosh\delta-\cos y)$, the periodised Cauchy kernel, so that $\pi\kappa_\delta=[P(y)-\ell_\delta(y)]-[P(0)-\ell_\delta(0)]$.

Case $y\ge\pi$. Here $P(y)\ge P(\pi)=\frac12\tanh(\delta/2)$ and $\ell_\delta(y)\le\ell_\delta(\pi)$, so
\[
\pi\kappa_\delta(y)\ge\tfrac12\tanh\tfrac\delta2-\tfrac12\coth\tfrac\delta2+\frac1\delta-\frac{\delta}{\delta^2+\pi^2}=\frac{\pi^2}{\delta(\delta^2+\pi^2)}-\frac{1}{\sinh\delta}>0,
\]
because $\pi^2\sinh\delta>\pi^2\delta+\pi^2\delta^3/6>\delta(\delta^2+\pi^2)$, using $\pi^2/6>1$.

Case $0<y<\pi$ and $\delta\le\sqrt3\,\pi$. Since $\ell_\delta''(x)=2\delta(3x^2-\delta^2)/(\delta^2+x^2)^3$, $\ell_\delta$ is strictly convex on $[\delta/\sqrt3,\infty)\supset[\pi,\infty)$, and $[2\pi m-y,2\pi m+y]\subset[\pi,\infty)$. So every bracket in \eqref{cg:eq:kappaU} is positive.

Case $0<y<\pi$ and $\delta\ge4.1$. Use $1-\cos y\le y^2/2$ and $\cosh\delta-\cos y\ge\cosh\delta-1=2\sinh^2(\delta/2)$ in the second term of \eqref{cg:eq:kappaU}, and $y^2/\{\delta(\delta^2+y^2)\}\ge y^2/\{\delta(\delta^2+\pi^2)\}$ in the first. Then $\pi\kappa_\delta>0$ provided $R(\delta)=8\tanh(\delta/2)\sinh^2(\delta/2)/\{\delta(\delta^2+\pi^2)\}>1$. The logarithmic derivative of the numerator exceeds $\coth(\delta/2)>1$, while that of the denominator, $1/\delta+2\delta/(\delta^2+\pi^2)$, is at most $2/\pi<1$ for $\delta\ge\pi$. So $R$ is increasing on $[\pi,\infty)$, and $R(4.1)=1.0322$ (computed, $1.032196804$). Since $4.1<\sqrt3\,\pi=5.4414$, the three cases cover all $\delta>0$ and $y>0$. The limit statements are Corollary~\ref{cg:cor:lorentz} together with $\varphi_U(s)=e^{-s/2}\sinh(s/2)/(s/2)$.
\end{proof}

The closed form agrees with the Lorentzian sum to fifteen digits at nine points ($\delta=0.1,1,5$, $y=0.5,3,7$), $\kappa_\delta>0$ on a grid of step $1/40$ on $(0,40]$ for ten values of $\delta$ from $0.01$ to $20$, and $R(4)=0.98037<1$, so the threshold $4.1$ in the proof is nearly sharp (computed). The mechanism is Poisson summation: sum over the integers minus integral is a sum of Lorentzians centred at the dual lattice $2\pi\mathbb{Z}$, and convexity of the Cauchy kernel makes them dominate.

\begin{theorem}[The shifted zeta as a Thorin Mellin transform]\label{cg:thm:mellinthorin}
For $1<\operatorname{Re}s<3$ and every $\delta>0$,
\begin{equation}\label{cg:eq:mellinthorin}
\zeta_\delta(s)=(2\pi)^s\int_0^\infty y^{-s}\kappa_\delta(y)\,dy,
\end{equation}
where $\zeta_\delta$ is the shifted family of Section~\ref{cg:sec:cmfamily}, continued analytically (for $\delta<2\pi$ it is the series \eqref{cg:eq:zetadelta}). At $\delta=0$ the right side is $(2\pi)^s\sum_m(2\pi m)^{-s}=\zeta(s)$.
\end{theorem}

\begin{proof}
Differentiate \eqref{cg:eq:thorinrep} in $w$: $h(w):=g(\delta+w)-g(\delta)=\int_0^\infty2w(w^2+y^2)^{-1}\kappa_\delta(y)\,dy$. Since $h(w)=O(w)$ at $0$ and $h(w)\to-g(\delta)$ at $\infty$, the integral $\int_0^\infty w^{-s}h(w)\,dw$ converges for $1<\operatorname{Re}s<2$. It is the analytic continuation of $\pi I_\delta(s)=\int w^{-s}g(\delta+w)\,dw$ from $0<\operatorname{Re}s<1$: subtracting the value at $0$ shifts the Mellin strip. By Tonelli ($\kappa_\delta\ge0$) and $\int_0^\infty w^{1-s}\,2(w^2+y^2)^{-1}dw=\pi y^{-s}/\sin(\pi s/2)$, we get $I_\delta(s)=\{\sin(\pi s/2)\}^{-1}\int y^{-s}\kappa_\delta(y)\,dy$. Multiply by $(2\pi)^s\sin(\pi s/2)$. The right side of \eqref{cg:eq:mellinthorin} converges for $1<\operatorname{Re}s<3$, since $\kappa_\delta(y)=O(y^2)$ at $0$ and $0\le\kappa_\delta\le1/(\pi\delta)$, which gives convergence for $\operatorname{Re}s>1$ and the continuation to that strip.
\end{proof}

At $s=2.5$, $\delta=0.5$ the right side of \eqref{cg:eq:mellinthorin} is $1.53343255$ (computed with $N=120$ periods of the lattice, error estimate $5\times10^{-8}$; $N=60$ gives $1.5334325$) and the series \eqref{cg:eq:zetadelta} gives $1.53343255193$ (computed). The theorem explains the identity of Comtet and Majumdar: formally \eqref{cg:eq:B1} is the Mellin transform of the Thorin measure $\sum_m\delta_{2\pi m}$ of the uniform law, divergent in the strip $0<s<1$, and the shift smooths this lattice into the positive density $\kappa_\delta$. This reading is for the shifted family; for the literal cutoff the Mellin transform is negative at real $s\in(1,2)$ (Remark~\ref{cg:rem:literal}).

\subsection{The Riemann function, by citation}\label{cg:sec:xi}

The same construction applies to $\xi$ and is in the book in full. The L\'evy measure of $\xi$ is primes minus a smooth density, $M_\xi=\sum_{n\ge2}\Lambda(n)\delta_{\log n}-e^x\,dx+dx/(e^{2x}-1)$, which is $-\sum_\rho e^{\rho x}$ as a distribution by the explicit formula (Proposition~\ref{bp:prop:levy}; $\Lambda(n)$ is von Mangoldt's function, which this chapter does not otherwise use). Its Thorin density at the basepoint $\alpha$ is $\pi^{-1}$ times the compensated density $v_\alpha(y)-b_\alpha$ of Chapter~\ref{chap:basepoint} (formula \eqref{bp:eq:excess}), and Proposition~\ref{bp:prop:valpha} proves, with no hypothesis on the zeros, that it is positive for every $\alpha\ge1$ and $y\ne0$, so the compensated reciprocal is a GGC for every $\alpha\ge1$; the threshold $7.789$, the root $7.789185$ of $\frac18\zeta(3,\alpha/2+1)=(\alpha-1)^{-3}$ (computed), also appears there. The unconditional GGC property of the uncompensated reciprocal for $\alpha\ge1$ is Theorem~\ref{thm:poisson}. Positivity of the uncompensated density for every $\alpha>\frac12$ is equivalent to RH by Lagarias's criterion (Corollary~\ref{pol:cor:equiv}), and under RH the passage to $\alpha\downarrow\frac12$ gives Hadamard's product exactly as Corollary~\ref{cg:cor:lorentz} does for the uniform law (Proposition~\ref{bp:prop:thorin}). Two levels of positivity are distinguished there: level (1), positivity of the uncompensated density, the GGC property of $\xi(\alpha)/\xi(\alpha+\sqrt s)$; and level (2), positivity of the compensated density, the GGC property of the compensated reciprocal $F_\alpha$, which strips the stable component $e^{-b\sqrt s}$ of index $\frac12$ out of level (1). Level (2) implies level (1); no basepoint is known at which (1) holds and (2) fails.

The uniform law is the model in which only level (2) holds. There the uncompensated density is $\pi^{-1}\operatorname{Re}g(\delta+iy)=\pi^{-1}[P(y)-\frac12-\ell_\delta(y)]$, which is negative near $y=\pi$ because $P(\pi)=\frac12\tanh(\delta/2)<\frac12$; the factor $e^{-s/2}$ in $\varphi_U$ is a drift, and compensation is what makes Theorem~\ref{cg:thm:uniform} work.

\begin{remark}[Zeta of the zeros]\label{cg:rem:zetazeros}
Under RH, as $\alpha\downarrow\frac12$ the Thorin density $\kappa^\xi_\alpha$ of $\xi$ (\eqref{cg:eq:kappa} with $f=\xi$) tends vaguely to $\sum_{\gamma>0}\delta_\gamma$, and the argument of Theorem~\ref{cg:thm:mellinthorin} suggests $\int_0^\infty y^{-s}\kappa^\xi_\alpha(y)\,dy\to\sum_{\gamma>0}\gamma^{-s}$, the zeta function of the zeros (Voros \cite{Voros10}), for $1<\operatorname{Re}s<3$. For the uniform law the same limit is $(2\pi)^{-s}\zeta(s)$ by Theorem~\ref{cg:thm:mellinthorin}. This is heuristic: the limit interchange needs the bounds $|\kappa^\xi_\alpha(y)|\le C\varepsilon y^2$ near the origin and a tail bound of order $\gamma^{-\sigma}+\varepsilon\gamma^{-2}$ for each Lorentzian against $y^{-\sigma}$, summed with the zero counting estimate, which are sketched and not written out. The range $\operatorname{Re}s>1$ of the source is wrong: the integral diverges at the origin for real $s\ge3$ because $\kappa^\xi_\alpha(y)=C_\varepsilon y^2+O(y^4)$ with $C_\varepsilon>0$ for small $\varepsilon$.
\end{remark}

\begin{table}[ht]
\centering
\small
\fitwidth{\begin{tabular}{p{2.9cm}p{5.4cm}p{5.4cm}}
\toprule
 & uniform law, shifted & Riemann $\xi$\\
\midrule
function & $\varphi_U(s)=(1-e^{-s})/s$ & $\xi(s)$\\
signed $k$ measure $M$ & $dx-d\lfloor x\rfloor$, Lebesgue minus integers & $\sum\Lambda(n)\delta_{\log n}-e^xdx+\frac{dx}{e^{2x}-1}$, primes minus smooth\\
duality & Poisson summation & explicit formula\\
zeros & $2\pi ik$, on $\operatorname{Re}s=0$ & $\rho$; on $\operatorname{Re}s=\tfrac12$ under RH\\
tilt & $\alpha=\delta$ (the shift) & $\alpha=\tfrac12+\varepsilon$\\
Thorin density & Lorentzian smoothing of $\sum\delta_{2\pi m}$, positive for all $\delta>0$ (Theorem~\ref{cg:thm:uniform}); only level (2) & uncompensated and compensated: positive for $\alpha\ge1$ (Theorem~\ref{thm:poisson}, Proposition~\ref{bp:prop:valpha}); uncompensated for all $\alpha>\tfrac12$ iff RH\\
limit Thorin measure & $\sum_m\delta_{2\pi m}$ & $\sum_{\gamma>0}\delta_\gamma$ counted with multiplicity, under RH\\
Mellin of Thorin & $(2\pi)^{-s}\zeta(s)$, proved & $\sum_{\gamma>0}\gamma^{-s}$, heuristic\\
\bottomrule
\end{tabular}}
\caption{The uniform law and the Riemann problem under the $2\sin^2$ lemma. Convergence of the Thorin densities is vague convergence. Without RH the limit at $\alpha=\frac12$ records only the zeros on the critical line (zeros at $\frac12\pm a+i\gamma$ cancel in the real part), so the limit measure is not a test of RH by itself.}
\label{cg:tab:parallel}
\end{table}
